\documentclass[11pt]{article}
\usepackage[margin=1in]{geometry}
\usepackage{amsmath,amssymb,amsthm}
\usepackage{graphicx}
\usepackage[hidelinks]{hyperref}

\newtheorem{theorem}{Theorem}
\newcommand{\RR}{\mathbb R}

\title{A Literal Counterexample to the Unit-Ball Extension Clause in Problem 22}
\author{Packet for arXiv:2602.15525}
\date{2026-06-21}

\begin{document}
\maketitle

\section*{Source Question}

The source is S.A. Bogaty and A.A. Tuzhilin,
\emph{Mazur--Ulam Theorem With Gromov--Hausdorff Distance},
arXiv:2602.15525. Problem 22 asks:

\begin{center}
\includegraphics[width=\textwidth]{figures/open_problem_crop.png}
\end{center}

In transcription, the problem asks whether, if Banach spaces \(V\) and \(W\)
admit an isometric embedding \(f:B_V\to W\), then there is a linear isometric
embedding \(F:V\to W\) such that \(F|_{B_V}=f\).

\section*{Answer}

The literal extension statement is false. There are Banach spaces \(V,W\) and
an isometric embedding \(f:B_V\to W\) which cannot be the restriction of any
linear map \(V\to W\), hence cannot be the restriction of a linear isometric
embedding.

\section*{Proof Intuition}

The obstruction is that an isometric embedding of a ball need not remember the
linear oddness of the ambient space. The map \(v\mapsto(v,\|v\|)\) places the
unit ball on the graph of the norm in an \(\ell_\infty\)-sum. Distances are
unchanged because the norm coordinate is 1-Lipschitz. But the graph is not
centrally symmetric: \(u\) and \(-u\) both have second coordinate \(1\). A
linear extension would have to send \(-u\) to the negative of the image of
\(u\), which is impossible.

\begin{theorem}
Let \(V\) be any nonzero real Banach space. Put
\[
  W=V\oplus_\infty \RR,\qquad
  \|(x,t)\|_W=\max\{\|x\|_V,|t|\}.
\]
Then
\[
  f:B_V\to W,\qquad f(v)=(v,\|v\|_V),
\]
is an isometric embedding, but no linear map \(F:V\to W\) satisfies
\(F|_{B_V}=f\).
\end{theorem}

\begin{proof}
For \(v,w\in B_V\),
\[
  \|f(v)-f(w)\|_W
  =\max\{\|v-w\|_V,\ |\|v\|_V-\|w\|_V|\}.
\]
By the reverse triangle inequality,
\[
  |\|v\|_V-\|w\|_V|\le \|v-w\|_V.
\]
Thus \(\|f(v)-f(w)\|_W=\|v-w\|_V\), so \(f\) is an isometric embedding of
\(B_V\) into \(W\).

Choose \(u\in V\) with \(\|u\|_V=1\). If a linear map \(F:V\to W\) satisfied
\(F|_{B_V}=f\), then
\[
  F(u)=f(u)=(u,1)
  \quad\text{and}\quad
  F(-u)=f(-u)=(-u,1).
\]
Linearity also gives
\[
  F(-u)=-F(u)=(-u,-1),
\]
which contradicts \(F(-u)=(-u,1)\). Therefore no such linear extension exists.
\end{proof}

\section*{Concrete One-Dimensional Instance}

Taking \(V=\RR\) and \(W=\ell_\infty^2\), the map
\[
  [-1,1]\ni t\longmapsto(t,|t|)
\]
is an isometric embedding, since
\[
  \max\{|s-t|,\ ||s|-|t||\}=|s-t|
\]
for all \(s,t\in[-1,1]\). It cannot be the restriction of a linear map
\(\RR\to\ell_\infty^2\), because it sends both \(1\) and \(-1\) to points
with second coordinate \(1\).

\section*{Scope and Limitations}

This packet refutes the literal parenthetical extension requirement
\(F|_{B_V}=f\). It does not refute the weaker reading that asks only whether
\(W\) contains some linear isometric copy of \(V\) whenever \(B_V\) embeds
isometrically into \(W\). In the construction above, the weaker conclusion is
true, because \(v\mapsto(v,0)\) is a linear isometric embedding of \(V\) into
\(V\oplus_\infty\RR\).

The same source paper recalls closely related nonlinear norm-graph
constructions before stating Problem 22. Therefore the human review question is
partly editorial: this should be counted either as a counterexample to Problem
22 as written, or as an internal wording issue rather than as a new answer to
the weaker intended problem.

\section*{Verification and Novelty Check}

The proof uses only the reverse triangle inequality and linearity. The included
scalar sanity script checks the one-dimensional instance on a finite grid; it
is not used as proof.

Local run indexes were searched for arXiv:2602.15525, \texttt{pr:boll},
\texttt{unit ball embedding}, and exact phrases involving \(F|_{B_V}=f\). Web
searches on 2026-06-21 for exact variants of Problem 22 found the source arXiv
record and no separate later answer. The novelty confidence is intentionally
limited because the counterexample is a one-line consequence of a construction
already adjacent in the source paper.

\begin{thebibliography}{9}
\bibitem{BT}
S.A. Bogaty and A.A. Tuzhilin,
\emph{Mazur--Ulam Theorem With Gromov--Hausdorff Distance},
arXiv:2602.15525, 2026.
\end{thebibliography}

\end{document}
